% !TEX root = ../main.tex
\chapter{Metric Manifolds}\label{ch:c10}
\begin{paracol}{2}

\begin{sources}
\begin{tabularx}{\linewidth}{@{}l l L@{}}
\toprule
\textbf{Source} & \textbf{Length} & \textbf{Content}\\
\midrule
\seg{10}{1}{Lecture 10: Metric Manifolds} & 1:20:28 & The metric as a symmetric, non-degenerate $(0,2)$ tensor field; the musical map and the so-called inverse metric; the round sphere of radius $R$; signature (Riemannian, Lorentzian, pseudo-Riemannian); speed and length of a curve and the length of the equator with a funny parametrisation; reparametrisation invariance; geodesics as stationary curves and the Euler--Lagrange equations of $L=\sqrt{g(\dot x,\dot x)}$; the Christoffel symbols and the Levi-Civita connection; the Riemann--Christoffel, Ricci, scalar and Einstein curvatures. Student questions on the inverse, units, compatibility and the trace of $G$.\\
\seg{10}{2}{Tutorial 9 \& 10: Metric Manifolds} & 1:17:31 & Solutions of the sheet by a tutor (not Schuller): true/false statements on the Levi-Civita connection, metrics and geodesics; signatures from null vectors; uniqueness of the Levi-Civita connection; reparametrisation and the square root in the Lagrangian; the geodesic equation of the round sphere; the Riemann--Christoffel tensor in normal coordinates and its symmetries.\\
\bottomrule
\end{tabularx}

\smallskip
\textbf{Transcripts} \file{materials/transcripts/C10-V01.txt} and \file{C10-V02.txt} (YouTube captions).
\textbf{Sheets 9 and 10:} see \url{https://tales.mbivert.com/on-heraeus-winter-school-tutorials/}.
\end{sources}

\begin{reading}[Where this chapter sits in the other notes]
\emph{Geometric Anatomy of Theoretical Physics} notes: the metric manifolds, the Levi-Civita connection and the curvature tensors follow the connections of chapters 21 to 25 there [exact chapter numbers and titles not checked]. In this book the connection and its coefficient functions are chapter~\ref{ch:c07}, curvature and torsion are chapter~\ref{ch:c08} (not yet written), and chapter~\ref{ch:c09} used the Ricci tensor as a result. The tangent bundle, in which the Lagrangian below lives, is chapter~\ref{ch:c06}; the symmetric $(0,2)$ tensors are chapter~\ref{ch:c03}.
\end{reading}

% ==================================================================
\section{The idea of a metric}\label{sec:c10-idea}

Welcome to the tenth lecture \ts{10}{1}{0:00}. ``We now return to mathematics proper'': the excursion to Newtonian physics (chapter~\ref{ch:c09}) is completed, and one again thinks abstractly and purely mathematically, pushing further the structures needed for general relativity. Before, a covariant derivative was established; now a so-called \textbf{metric}. The idea is a structure on a smooth manifold that allows to assign to vectors in each tangent space a length; a norm would do for that, but one is more special and wants also an \emph{angle} between vectors (of the same tangent space): a geometric structure on a vector space giving both length and angle is an \textbf{inner product} (sometimes ``scalar product''). One wants an inner product on \emph{each} tangent space, an ``inner product field''; and an inner product is a $(0,2)$ tensor (chapter~\ref{ch:c03}) with certain properties. That is a metric
\ts{10}{1}{0:21}\ts{10}{1}{0:41}\ts{10}{1}{1:02}\ts{10}{1}{1:25}\ts{10}{1}{1:46}\ts{10}{1}{2:06}\ts{10}{1}{2:27}.

\paragraph{The logic.} From the metric one then defines the length of a curve: if one knows how to turn the velocity vector into a number, the \emph{speed}, one can integrate it over the parameter and get the distance. So the order is the other way round from school: velocity first, then speed, then the length of a curve (``usually we think the length of a curve is clear and the velocity more complicated''). With the length one can look at shortest curves, also longest curves, and curves whose length does not change if one moves the curve a little, the \emph{stationary} curves; shortest ones are called \textbf{geodesics}.
In flat space the shortest curve is also the straightest, and one already knows how to speak about straightest curves with a covariant derivative (auto-parallels). If one establishes the metric and the covariant derivative fully independently, the shortest curves are of course not the straightest ones: ``you could get anything''. But one can \emph{require} that the shortest curves be the straightest curves and the other way round: this is a compatibility criterion (\emph{straight $=$ short}, more precisely straight $=$ stationary), which makes a \emph{metric compatible} connection, and the connection will be fully determined by the metric (in the presence of zero torsion). From a metric one then gets a connection, from it a curvature, and so on: that is the outline
\ts{10}{1}{2:48}\ts{10}{1}{3:09}\ts{10}{1}{3:29}\ts{10}{1}{3:51}\ts{10}{1}{4:13}\ts{10}{1}{4:33}\ts{10}{1}{4:54}\ts{10}{1}{5:16}\ts{10}{1}{5:37}\ts{10}{1}{5:58}.

\section{Metrics}\label{sec:c10-metric}

\begin{definition}[New concepts: metric, musical map]\label{def:c10-metric}
A \textbf{metric} $g$ on a smooth manifold $(M,\mathcal O,\mathcal A)$ (at first there is no covariant derivative; this is structure thrown in by hand) is a $(0,2)$ tensor field satisfying two conditions:
\begin{enumerate}
  \item \emph{Symmetry}: for all vector fields $X,Y$,
        \[g(X,Y)=g(Y,X).\]
  \item \emph{Non-degeneracy}: the \textbf{musical map} ``flat'' (as in music),
        \[\flat\colon\Gamma(TM)\to\Gamma(T^*M),\qquad X\mapsto X^\flat,\qquad X^\flat(Y):=g(X,Y),\]
        which takes a vector field $X$ to the covector field obtained by keeping one slot of $g$ open (the covector field can act on a vector field $Y$ and produces a function, which is declared to be $g(X,Y)$), is a $C^\infty(M)$-\emph{isomorphism}, i.e.\ invertible.
\end{enumerate}
\end{definition}
\ts{10}{1}{6:19}\ts{10}{1}{6:40}\ts{10}{1}{7:01}\ts{10}{1}{7:21}\ts{10}{1}{7:42}\ts{10}{1}{8:04}\ts{10}{1}{8:24}\ts{10}{1}{8:45}\ts{10}{1}{9:06}\ts{10}{1}{9:27}\ts{10}{1}{9:48}\ts{10}{1}{10:09}\ts{10}{1}{10:30}.
(A metric is thus a $(0,2)$ tensor field that is ``imposed by hand''. Schuller says later that what he defines might justifiedly be called a \emph{pseudo}-metric by others, since no signature is specified; only the invertibility of the musical map is required, see below.) The musical map is the new thing: by virtue of the metric one has a means of converting a vector field into a covector field (``in quantum mechanics you could say you send a ket to a bra''); in components $(X^\flat)_a=g_{am}X^m$.

\paragraph{Lowering and raising indices.} Some people say: if the flat map confuses, do not write it, just put the index down; the covector that comes with the vector. This is called lowering the index with the metric. The inverse of the musical isomorphism (called ``sharp'' by music enthusiasts; ``we do not'') takes a covector field to a vector field, carrying an index up: raising the index. In components it all looks particularly simple. ``There is a nice statement about this whole technology of pulling indices up and down. It is not needed. Nobody needs this and we will not do it'': the musical map plays a role because it \emph{must exist and be invertible}, that is the non-degeneracy condition
\ts{10}{1}{10:54}\ts{10}{1}{11:15}\ts{10}{1}{11:36}\ts{10}{1}{11:57}\ts{10}{1}{12:19}\ts{10}{1}{12:39}\ts{10}{1}{12:59}\ts{10}{1}{13:20}.

\begin{definition}[New concepts: the so-called inverse metric]\label{def:c10-inverse}
The \textbf{so-called inverse metric} $g^{-1}$ of a metric $g$ is the symmetric $(2,0)$ tensor field which takes two covector fields $\omega,\sigma$ to the smooth function
\[
  g^{-1}(\omega,\sigma):=\omega\bigl(\sigma^\sharp\bigr)=\omega\bigl(\flat^{-1}(\sigma)\bigr),
\]
that is: first one raises the index of $\sigma$ by the inverse musical isomorphism, getting a vector field, and then one lets $\omega$ act on it. In a chart its components $g^{ab}$ satisfy $g^{am}g_{mb}=\delta^a{}_b$ (summed over $m$).
\end{definition}
\ts{10}{1}{13:44}\ts{10}{1}{14:08}\ts{10}{1}{14:29}\ts{10}{1}{14:50}\ts{10}{1}{15:10}\ts{10}{1}{15:31}\ts{10}{1}{15:51}.
(Schuller first used the inverse musical map and said ``inverse'' too early: ``I jumped ahead, it should have come first''.) Why ``so-called''? The metric is a map $\Gamma(TM)\times\Gamma(TM)\to C^\infty(M)$, and its inverse would have to go the other way round; $g^{-1}$ is a map on covector fields, as written. It is an inverse in the sense of the nice condition $g^{am}g_{mb}=\delta^a{}_b$, and if one thinks of matrices and forgets everything about the valence of tensors one can say it is an inverse matrix; ``the logic goes via this''
\ts{10}{1}{16:12}\ts{10}{1}{16:33}.
\begin{coursenote}[Convention on raised indices]
Whenever one writes $g$ with upstairs indices, which makes no sense for a $(0,2)$ tensor field, one means this so-called inverse metric, a $(2,0)$ tensor field; ``it is a totally different tensor, it just has the indices up by convention'', these indices are not pulled up by anything. Schuller never suppresses the metric in a formula, hence he does not use the pull-up/pull-down technology: an object whose components have an index down could be an original covector or one constructed from a vector by a metric; ``I do not want to forget which is which. We will never suppress the metric in our formulae, never ever''
\ts{10}{1}{16:54}\ts{10}{1}{17:14}\ts{10}{1}{1:17:24}\ts{10}{1}{1:17:44}\ts{10}{1}{1:18:05}.
\end{coursenote}

\begin{example}[The round sphere of radius $R$]\label{ex:c10-sphere}
Take the smooth manifold $S^2$ with the topology and smooth atlas established before and the chart $(U,x)$ with $x=(\theta,\varphi)$ (the two angles), which covers most of the sphere apart from a half equator line that has to be taken out ($\varphi\in(0,2\pi)$, $\theta\in(0,\pi)$): ``a pretty almost global chart''. The smooth sphere has \emph{no} metric (``it can be a potato shape-wise, it has no shape''), also no covariant derivative. Define the metric by its components, evaluated at the point described by the coordinates in the chart,
\[
  g_{ij}=\begin{pmatrix}R^2&0\\0&R^2\sin^2\theta\end{pmatrix},\qquad R\in\R^+\text{ a constant},
\]
(rows and columns just a convention to say ``look at this row and this column''; ``do not think this is a matrix that can be multiplied''). This defines the metric of the \emph{round sphere of radius $R$}: before, the round sphere as a manifold with a connection had a shape but no size; ``the round sphere as a manifold with metric has everything fixed'', what one always wanted fixed. On the same smooth space one could establish a different metric and make it an ellipsoid of axes $A,B,C$, or a potato: whatever metric one induces gives the space \emph{shape and size at the same time} (intuitively, ``we will come to that'').
\end{example}
\ts{10}{1}{17:56}\ts{10}{1}{18:17}\ts{10}{1}{18:37}\ts{10}{1}{18:58}\ts{10}{1}{19:19}\ts{10}{1}{19:39}\ts{10}{1}{20:00}\ts{10}{1}{20:21}\ts{10}{1}{20:42}.
This metric is defined in one chart only; to define it globally one must prescribe the component functions in another chart as well, and make sure that in the overlap the component functions with respect to one chart transform into those of the other according to the \emph{tensor field component transformation law} (as a $(0,2)$ tensor field), otherwise inconsistent choices have been made \ts{10}{1}{21:02}\ts{10}{1}{21:23}\ts{10}{1}{21:44}.

\section{Signature}\label{sec:c10-signature}

``A topic from linear algebra.'' For a $(1,1)$ tensor with components $a^a{}_m$, the eigenvalue equation $a^a{}_mv^m=\lambda v^a$ makes sense (summed index up and down; the outcome is a vector). For a $(0,2)$ tensor with components $g_{am}$ one can multiply a vector in, $g_{am}v^m$, but one cannot write $=\lambda v_a$ with the same $v$: the left is a covector, the right a vector component. \emph{An eigenvalue is a good notion for a $(1,1)$ tensor and not for a $(0,2)$ tensor}; and ``symmetric'' for a $(1,1)$ tensor can only mean symmetric with respect to a $(0,2)$ tensor. The tensor of inertia of classical mechanics must be a $(1,1)$ tensor if written properly, otherwise one could not capture the different inertia scalars in the various directions
\ts{10}{1}{22:04}\ts{10}{1}{22:24}\ts{10}{1}{22:44}\ts{10}{1}{23:05}\ts{10}{1}{23:26}\ts{10}{1}{23:46}\ts{10}{1}{25:09}.
For a symmetric endomorphism (a symmetric $(1,1)$ tensor with respect to a metric) one finds a basis in which it is diagonal: ``it is often said, but still wrong, that you can also bring this to eigenvalue form'' in this sense: for a symmetric bilinear form one can bring it to a form consisting of $+1$'s, $-1$'s and $0$'s, but the numbers are \emph{not} eigenvalues, ``you will never find a basis that will scale away whatever these $\lambda$'s are, but you will always find a basis in which whatever you have can be scaled to these numbers'' \ts{10}{1}{23:46}\ts{10}{1}{24:08}\ts{10}{1}{24:28}\ts{10}{1}{24:48}.

\begin{definition}[New concepts: signature, Riemannian, Lorentzian, pseudo-Riemannian]\label{def:c10-signature}
A symmetric $(0,2)$ tensor at a point (on a vector space $V$) has a \textbf{signature} $(p,q)$: there is a basis in which it takes the diagonal form with $p$ entries $+1$, $q$ entries $-1$ and $\dim V-p-q$ entries $0$. The numbers $p,q$ are well defined (``you can change bases but you will never change the number of pluses and minuses''): this is much less information than $\dim V$ real eigenvalues, it is a pair of non-negative integers, and it is coordinate independent. For a non-degenerate one (musical map invertible: ``no zeros'') $p+q=\dim V$, and there are $\dim V+1$ possible signatures, e.g.\ for $\dim V=3$: $(+,+,+)$, $(+,+,-)$, $(+,-,-)$, $(-,-,-)$ (the others are reorderings). On a manifold a metric $g$ is called
\begin{itemize}
  \item \textbf{Riemannian} if its signature is $(+,\dots,+)$ (an overall minus sign, only minuses, would in a sense also be Riemannian, but most people say plus);
  \item \textbf{Lorentzian} if its signature is $(+,-,\dots,-)$, i.e.\ one plus and the others minus ($1$ plus and $3$ minuses in four dimensions; an overall sign does not change the game but a convention must be fixed, and this will usually be the one used here);
  \item \textbf{pseudo-Riemannian} for all other signatures, the Lorentzian ones included.
\end{itemize}
\end{definition}
\ts{10}{1}{25:29}\ts{10}{1}{25:50}\ts{10}{1}{26:10}\ts{10}{1}{26:31}\ts{10}{1}{26:52}\ts{10}{1}{27:12}\ts{10}{1}{27:33}\ts{10}{1}{27:53}\ts{10}{1}{28:16}\ts{10}{1}{28:38}\ts{10}{1}{28:58}\ts{10}{1}{29:19}.
Only the Riemannian metrics define in each tangent space an \emph{inner product}, one property of which is $g(X,X)=0$ if and only if $X=0$: every non-zero vector has non-zero length; this is not the case for the others. For relativity the Lorentzian one will be the one wanted on spacetime, but ``we are not there yet''; when Schuller says metric, others might justifiedly call it a pseudo-metric, because he did not specify the signature. The signature is a property of the $(0,2)$ tensor that is independent of the choice of chart; to display it one needs a chart, and at one point one can bring $g$ to diagonal form by choice of chart, but at the next point the entries may all be filled in, and the inner product changes from point to point, while the signature is the same ``and the object carries it''. In the Lorentzian case, in each tangent space $T_pM$ there are vectors that are \emph{null}, $g(X,X)=0$ although non-zero: they give the cones (``go to the tutorials, there you will learn what these different signatures have to do with cones in the tangent spaces; it is very important'') \ts{10}{1}{29:19}\ts{10}{1}{29:39}\ts{10}{1}{30:00}\ts{10}{1}{30:21}\ts{10}{1}{30:41}\ts{10}{1}{32:08}\ts{10}{1}{32:29}.

\section{Speed and length of a curve}\label{sec:c10-length}

Why afford such a structure? In order to define the length of curves; it is all right to measure length and angles in a tangent space, but is that all one can do? Let $\gamma$ be a smooth curve; at each $\gamma(\lambda)$ its velocity is known, but not its speed: ``the speed we cannot define on a smooth manifold''. On a metric manifold we can; to avoid too many case distinctions let it be a Riemannian metric manifold \ts{10}{1}{30:41}\ts{10}{1}{31:04}\ts{10}{1}{31:27}\ts{10}{1}{31:48}\ts{10}{1}{32:49}.

\begin{definition}[New concepts: speed, length of a curve]\label{def:c10-speed}
On a Riemannian metric manifold $(M,\mathcal O,\mathcal A,g)$ the \textbf{speed} of a smooth curve $\gamma$ at $\gamma(\lambda)$ is the number
\[
  \sqrt{g\bigl(v_\gamma,v_\gamma\bigr)\big|_{\gamma(\lambda)}},
\]
where $v_\gamma$ is the velocity vector field of the curve (plugged in twice, hence the square root). For a smooth curve $\gamma\colon[0,1]\to M$ (``one can always do this with an interval'') the \textbf{length} of $\gamma$ is the real number (a \emph{functional}, ``an action functional''; it eats a curve and yields a number)
\[
  L[\gamma]:=\int_0^1\mathrm d\lambda\,\sqrt{g\bigl(v_\gamma,v_\gamma\bigr)\big|_{\gamma(\lambda)}}.
\]
\end{definition}
\ts{10}{1}{32:49}\ts{10}{1}{33:10}\ts{10}{1}{33:31}\ts{10}{1}{33:51}\ts{10}{1}{38:24}\ts{10}{1}{38:45}\ts{10}{1}{39:06}\ts{10}{1}{39:27}\ts{10}{1}{39:47}.
So velocity is more fundamental than speed, and speed more fundamental than the length of a curve: ``exactly the opposite way you would go in high school or kindergarten''.
\begin{aside}[The need for speed]
``I feel the need, the need for speed'': Maverick in \emph{Top Gun}; ``who knows that? Nobody. Okay, almost'': one who feels the need for speed feels the need for a metric. Velocity is defined, speed is not \ts{10}{1}{34:13}\ts{10}{1}{34:34}.
\end{aside}

\paragraph{Units.} The physical dimension of the velocity vector components is $[\dot x^a]=1/\text{time}$, ``how can that be? Did we not learn that the components of the velocity are speeds in the different directions?'' No, because these components refer to an arbitrary chart, in which one may make big steps if one increases a coordinate by one, and smaller in another (bigger or smaller as measured by the metric, which one did not have yet). The components of the metric in the same chart carry units $\text{length}^2$ (for each coordinate combination), so $g_{ab}\dot x^a\dot x^b$ has $\text{length}^2/\text{time}^2$, and the square root gives length per time, reconciling with the notion of speed. ``The idea that coordinate distance has anything to do with real distance is false; that is what took Einstein a long time to totally lose, the notion that going double as far in coordinates has anything to do with going double as far in reality.'' The trick: the metric transforms as a $(0,2)$ tensor, the velocity components as vectors, the transformations cancel and the result is invariant; the metric carries the information how to translate coordinate distances into real length. (A student asks about the bracket in $[\dot x]$: physicist notation for the physical dimension.) ``In mathematics you do not put this, but as a physicist it is important to see''
\ts{10}{1}{34:34}\ts{10}{1}{34:57}\ts{10}{1}{35:17}\ts{10}{1}{35:39}\ts{10}{1}{35:59}\ts{10}{1}{36:20}\ts{10}{1}{36:41}\ts{10}{1}{37:01}.
\begin{coursenote}[Student question: does a metric imply a connection?]
Not yet: one can throw in a connection of one's choosing; a metric and a covariant derivative are both described by hand, and ``what on earth would they have to do with each other?'' One might get the idea that the shortest curves with respect to the metric coincide with the straight curves of the connection, but that is an additional assumption; and if one makes it one can express $\nabla$ in terms of $g$ (see below) \ts{10}{1}{37:22}\ts{10}{1}{37:43}\ts{10}{1}{38:04}.
\end{coursenote}

\begin{example}[The equator with a funny parametrisation]\label{ex:c10-equator}
Take the round sphere of radius $R$ (\cref{ex:c10-sphere}) and its equator: the curve $\gamma$ with $\theta=\pi/2$, parametrised on $[0,1]$ as
\[
  \theta(\lambda)=\frac\pi2,\qquad\varphi(\lambda)=2\pi\lambda^3+42
\]
(a $+42$ to stay safely inside the chart, where $\varphi\in(0,2\pi)$ is understood modulo $2\pi$ [?]; ``one point does not matter for the integral''); ``be cheeky: any crazy thing, as long as from the beginning to the end it makes the curve''. In the chart the velocity components are the derivatives of the chart representatives: $\theta'=0$, $\varphi'=6\pi\lambda^2$, so
\begin{align*}
  L[\gamma]&=\int_0^1\mathrm d\lambda\sqrt{R^2\,\theta'^2+R^2\sin^2\theta\,\varphi'^2}\\
  &=\int_0^1\mathrm d\lambda\,\sqrt{R^2\cdot1\cdot36\pi^2\lambda^4}\\
  &=\int_0^1 6\pi R\,\lambda^2\,\mathrm d\lambda=2\pi R,
\end{align*}
the circumference of the equator (the cross term vanishes as the metric has zeros off the diagonal; $\sin(\pi/2)=1$)
\end{example}
\ts{10}{1}{40:28}\ts{10}{1}{40:50}\ts{10}{1}{41:11}\ts{10}{1}{41:32}\ts{10}{1}{41:53}\ts{10}{1}{42:13}\ts{10}{1}{42:33}\ts{10}{1}{42:54}\ts{10}{1}{43:14}\ts{10}{1}{43:35}\ts{10}{1}{43:56}\ts{10}{1}{44:17}\ts{10}{1}{44:38}\ts{10}{1}{44:58}\ts{10}{1}{45:19}\ts{10}{1}{45:39}\ts{10}{1}{45:59}\ts{10}{1}{46:20}\ts{10}{1}{46:42}\ts{10}{1}{47:02}\ts{10}{1}{47:23}.
\begin{coursenote}[Slips at the board in the integral]
The lecture takes the parametrisation with $\lambda^3$ and, in the integrand, first writes $\varphi'=2\pi\cdot3\lambda^2$ and then forgets the square of the factor (``I am lacking this $36$; it is this one squared''): with $\varphi'^2=36\pi^2\lambda^4$ the integrand is $6\pi R\lambda^2$ and the integral gives $2\pi R$, as it must. ``The funny parametrisation gave us some funny factors, but in the end it all worked out''
\ts{10}{1}{45:39}\ts{10}{1}{45:59}\ts{10}{1}{46:20}\ts{10}{1}{46:42}\ts{10}{1}{47:02}.
\end{coursenote}

\begin{theorem}[Reparametrisation invariance of the length]\label{thm:c10-reparam}
Let $\gamma\colon[0,1]\to M$ be a smooth curve and $\sigma\colon[0,1]\to[0,1]$ smooth, bijective and \emph{increasing}. Then $L[\gamma\circ\sigma]=L[\gamma]$.
\end{theorem}
``Increasing'' because otherwise the curve could run backwards: ``if you drive and then your wife has forgotten something, you drive home again and then you drive on, then the length is longer; that is why when you go on holiday you only do increasing parametrisations of our parameter, and bijective'' (``the female lecturers will make the opposite jokes, so there will be equality''). The parametrisation certainly should not matter, and it does not; the proof is in the tutorial (\cref{sec:t10-1}, exercise 4) \ts{10}{1}{47:44}\ts{10}{1}{48:04}\ts{10}{1}{48:25}\ts{10}{1}{48:51}\ts{10}{1}{49:12}\ts{10}{1}{49:33}.

\section{Geodesics}\label{sec:c10-geodesics}

The interesting thing happens for extremal curves.

\begin{definition}[New concepts: geodesic]\label{def:c10-geodesic}
A curve $\gamma\colon[0,1]\to M$ on a Riemannian manifold $(M,\mathcal O,\mathcal A,g)$ is called a \textbf{geodesic} if it is a \emph{stationary curve} for the length functional $L$.
\end{definition}
\ts{10}{1}{49:33}\ts{10}{1}{49:55}\ts{10}{1}{50:19}\ts{10}{1}{50:40}.
(Later one can drop the Riemannian assumption and do it somewhat differently.) What is a stationary curve? ``You know this idea from classical mechanics'': one associates with the curve $\gamma$ a number $L[\gamma]$ (in classical mechanics the action, unfortunately not the $L$ of the Lagrangian: it is more like an $S$), and one can ask for deforming the curve a little, $\gamma\to\gamma+\epsilon\,\delta\gamma$ (``let us not write it formally, you look at the lectures of classical mechanics''): the length of the deformed curve agrees with $L[\gamma]$ to \emph{first order} in $\epsilon$, i.e.\ the derivative with respect to $\epsilon$ at $\epsilon=0$ vanishes; second order may differ.
One is tempted to say ``minimal'', but depending on the signature and where one is, a stationary curve could be a maximum: on the sphere the great-circle arc that is the longest between two points is also a geodesic (``the longest curve between these two points; this is the shortest; this is also a geodesic''), and it could even be a point of inflection of the length with respect to $\epsilon$. Hence not ``extremal'': \emph{stationary}
\ts{10}{1}{50:40}\ts{10}{1}{51:00}\ts{10}{1}{51:20}\ts{10}{1}{51:41}\ts{10}{1}{52:01}\ts{10}{1}{52:21}\ts{10}{1}{52:41}.

\begin{theorem}[Geodesics and the Euler--Lagrange equations]\label{thm:c10-el}
A curve $\gamma$ is a geodesic if and only if it satisfies the Euler--Lagrange equations of the Lagrangian $L\colon TM\to\R$,
\[
  L(x,v)=\sqrt{g(v,v)}\Big|_x,
\]
a function on the tangent bundle (``it eats a position and a velocity, i.e.\ a point of the tangent bundle''; purely geometrical, so without indices), which is just a mathematical reformulation of the stationarity condition.
\end{theorem}
\ts{10}{1}{53:04}\ts{10}{1}{53:24}\ts{10}{1}{53:45}.
In a chart the Euler--Lagrange equations take the form $\dfrac{\mathrm d}{\mathrm d\lambda}\dfrac{\partial L}{\partial\dot\gamma^m}-\dfrac{\partial L}{\partial\gamma^m}=0$: ``this is really a chart-dependent notation and you need to know what you mean by it''; the precise formulation uses the chart $\xi_x$ that was constructed on the tangent bundle (chapter~\ref{ch:c06}, $\dim TM=2\dim M$), the function $L$ on $TM$ and the \emph{lift} $\sigma$ of the curve to $TM$, $\sigma(\lambda)=(\gamma(\lambda),v_\gamma(\lambda))$: with the symbols of the charts of $TM$, $\dfrac{\mathrm d}{\mathrm d\lambda}\Bigl(\dfrac{\partial L}{\partial v^m}\Bigr)\bigl(\sigma(\lambda)\bigr)-\dfrac{\partial L}{\partial x^m}\bigl(\sigma(\lambda)\bigr)=0$.
Schuller does not derive the equations from the abstract formalism (``I would have needed another half an hour because you are not familiar with the formalism'') but uses the chart representative $L=\sqrt{g_{ij}(\gamma)\,\dot\gamma^i\dot\gamma^j}$, ``what you would write down in classical mechanics, apart from the square root, which is important to make it reparametrisation invariant: in classical mechanics you do not want that because you want the uniform parameter, that is why in classical mechanics you do not do the square root; but we are not doing classical mechanics, we are doing geometry''
\ts{10}{1}{54:06}\ts{10}{1}{54:28}\ts{10}{1}{54:50}\ts{10}{1}{55:10}\ts{10}{1}{55:31}\ts{10}{1}{55:51}\ts{10}{1}{1:06:16}\ts{10}{1}{1:06:36}\ts{10}{1}{1:06:57}.
\begin{coursenote}[Student question: derive with respect to $\dot\gamma$ as independent?]
The student's doubt is how one can derive with respect to $\dot\gamma$ and treat it as independent of $\gamma$. They are \emph{not} independent along the curve, but this has nothing to do with it: $L$ is a function on the tangent bundle, in which position and velocity are independent coordinates, and the curve is only plugged in \emph{after} the derivatives were taken; ``it is a little bit deviating from the geometric picture, but it is the right thing, of course'' \ts{10}{1}{1:07:17}\ts{10}{1}{1:07:38}\ts{10}{1}{1:07:58}.
\end{coursenote}

\paragraph{The computation.} Write the Euler--Lagrange equations of $L=\sqrt{g_{ij}(\gamma)\dot\gamma^i\dot\gamma^j}$. First
\[
  \frac{\partial L}{\partial\dot\gamma^m}=\frac{1}{2\sqrt{g_{ij}\dot\gamma^i\dot\gamma^j}}\,2g_{mj}\dot\gamma^j=\frac{g_{mj}\dot\gamma^j}{\sqrt{g_{ij}\dot\gamma^i\dot\gamma^j}}
\]
(the two from the symmetry cancels the two from the root). Take its derivative with respect to the parameter: by the product rule it involves the derivative of $1/\sqrt{\cdots}$ and then of $g_{mj}\dot\gamma^j$, and the latter contains $g_{mj}\ddot\gamma^j$ plus, because $g_{mj}$ depends on the point of the curve, the chain-rule term $\partial_sg_{mj}\,\dot\gamma^s\dot\gamma^j$. ``This term you mostly do not write down in classical mechanics, because there the metric is the diagonal $(1,1,1)$ that does not depend on where you are; but this term gives you a lot of insight about geometry.''
Because of \cref{thm:c10-reparam} one may reparametrise without changing the length, and chooses the parameter such that the speed is \emph{one} at every point (``I can drive from here to my home, if I want, at one meter per second, however the path is''; always possible). Then $g_{ij}\dot\gamma^i\dot\gamma^j=1$ for any $\lambda$, its derivative is zero and the derivative of the $1/\sqrt{\cdots}$ factor drops out. The right-hand part is
\[
  \frac{\partial L}{\partial\gamma^m}=\frac{1}{2\sqrt{\cdots}}\,\partial_mg_{ij}\,\dot\gamma^i\dot\gamma^j=\tfrac12\,\partial_mg_{ij}\,\dot\gamma^i\dot\gamma^j
\]
(``not zero although I do not have a potential: usually this term is only non-zero if you have a potential, but $g$ depends on the point directly''). Putting it together:
\[
  g_{mj}\ddot\gamma^j+\partial_sg_{mj}\,\dot\gamma^s\dot\gamma^j-\tfrac12\,\partial_mg_{ij}\,\dot\gamma^i\dot\gamma^j=0.
\]
Multiply by the so-called inverse metric $g^{qm}$ and sum over $m$ ($g^{qm}g_{mj}=\delta^q{}_j$); rename $s\to i$ in the second term (no conflicting index) and use that $\dot\gamma^i\dot\gamma^j$ is symmetric to symmetrise the bracket (``the whole equation in standard form, with a $\tfrac12$ in front of everything because I doubled this term''):
\begin{equation}\label{eq:c10-geodesic}
  \ddot\gamma^q+\tfrac12\,g^{qm}\bigl(\partial_ig_{mj}+\partial_jg_{mi}-\partial_mg_{ij}\bigr)\,\dot\gamma^i\dot\gamma^j=0.
\end{equation}
This is the \textbf{geodesic equation} for the components of $\gamma$ in a chart
\ts{10}{1}{55:51}\ts{10}{1}{56:13}\ts{10}{1}{56:33}\ts{10}{1}{56:54}\ts{10}{1}{57:14}\ts{10}{1}{57:34}\ts{10}{1}{57:56}\ts{10}{1}{58:16}\ts{10}{1}{58:38}\ts{10}{1}{58:59}\ts{10}{1}{59:20}\ts{10}{1}{59:40}\ts{10}{1}{1:00:01}\ts{10}{1}{1:00:22}\ts{10}{1}{1:00:42}\ts{10}{1}{1:01:03}\ts{10}{1}{1:01:25}\ts{10}{1}{1:01:45}\ts{10}{1}{1:02:06}\ts{10}{1}{1:02:27}\ts{10}{1}{1:02:48}\ts{10}{1}{1:03:10}\ts{10}{1}{1:03:30}\ts{10}{1}{1:03:51}\ts{10}{1}{1:04:12}\ts{10}{1}{1:04:33}\ts{10}{1}{1:04:54}.
\begin{coursenote}[A small slip in the board derivation]
In the lecture the index placement on the board is slightly off in places (the derivative index written once as $\partial/\partial x^i$ after a student's remark: ``so actually this should be $\partial/\partial x^i$''); the final equation above is the standard one \ts{10}{1}{1:05:35}.
\end{coursenote}

\paragraph{The Levi-Civita connection.} From the Newtonian chapter one observes at once: the whole combination in \cref{eq:c10-geodesic}, with an index $q$ up and $i,j$ down, depending only on where the curve passes, is an object that can be called a $\Gamma$, and \cref{eq:c10-geodesic} is an \emph{auto-parallel equation}:
\begin{definition}[New concepts: Christoffel symbols, Levi-Civita connection]\label{def:c10-lc}
The \textbf{Christoffel symbols} of the metric $g$ in a chart are
\[
  \Gamma^{(\mathrm{LC})\,q}{}_{ij}:=\tfrac12\,g^{qm}\bigl(\partial_ig_{mj}+\partial_jg_{mi}-\partial_mg_{ij}\bigr).
\]
They are the connection coefficient functions (\cref{def:c07-gamma}) of a connection $\nabla^{\mathrm{LC}}$, the \textbf{Levi-Civita connection} of $g$: the connection that identifies the geodesics with the auto-parallels. (The $\Gamma^{(\mathrm{LC})}$ themselves are also often called the Christoffel symbols, ``for traditional reasons'', although they are nothing but connection coefficient functions of a certain connection.)
\end{definition}
\ts{10}{1}{1:05:15}\ts{10}{1}{1:07:58}\ts{10}{1}{1:08:19}\ts{10}{1}{1:08:41}\ts{10}{1}{1:09:01}\ts{10}{1}{1:09:21}.
So insisting that the geodesic equation, which has nothing to do with auto-parallels, \emph{be also} the auto-parallel equation introduces a new structure through the $\Gamma$'s but one must choose them in a very specific way. This choice is usually made if a metric is given: one thinks of a smooth manifold with metric \emph{and} connection, the connection arising from the metric. Chart by chart it is this prescription; ``you can say this is chart dependent, does this really work? It does work'': an abstract way, to be treated in the tutorials, is that a connection $\nabla$ with $\nabla g=0$ and with vanishing torsion, $T^\nabla=0$, is already the Levi-Civita connection $\nabla^{\mathrm{LC}}$ (totally chart free). One can get it by requiring the shortest or stationary curves to be the auto-parallels, or without any reference to curves by $\nabla g=0$ everywhere and zero torsion: two points of view, the second ``very economical and practical'', called \textbf{metric compatibility}. What does it mean? That auto-parallels coincide with geodesics; and why do this? ``If you only have a metric you might want to think about the connection, because once you have that you can define curvature''
\ts{10}{1}{1:09:43}\ts{10}{1}{1:10:03}\ts{10}{1}{1:10:23}\ts{10}{1}{1:10:45}\ts{10}{1}{1:11:05}\ts{10}{1}{1:11:25}\ts{10}{1}{1:11:46}.
\begin{coursenote}[Torsion and the exact statement]
(Not said in class.) The theorem to be proved in the tutorial is: on a metric manifold there is exactly one connection that is torsion free and satisfies $\nabla g=0$, with the Christoffel symbols above. The torsion itself is defined in the next lecture (chapter~\ref{ch:c08}; lecture 7 only called the antisymmetric part of $\Gamma$ torsion).
\end{coursenote}

\section{The curvatures of a metric manifold}\label{sec:c10-curvature}

The last definitions: curvature on a metric manifold. If one has only a metric manifold and speaks of ``the Riemann curvature'', what one always means is the one of the metric-induced Levi-Civita connection, because curvature was defined in terms of a connection (chapter~\ref{ch:c08}, not yet written), and with a metric one has the Levi-Civita connection. The following definitions use charts, ``something ugly: I can do this component free, but it is really not worth it'' \ts{10}{1}{1:11:46}\ts{10}{1}{1:12:12}\ts{10}{1}{1:13:33}\ts{10}{1}{1:13:55}\ts{10}{1}{1:14:17}.

\begin{definition}[New concepts: Riemann--Christoffel, Ricci, scalar and Einstein curvature]\label{def:c10-curv}
Let $(M,\mathcal O,\mathcal A,g)$ be a metric manifold and $R^a{}_{bcd}$ the Riemann curvature tensor of its Levi-Civita connection.
\begin{enumerate}
  \item The \textbf{Riemann--Christoffel curvature} is the $(0,4)$ tensor field with components $R_{abcd}:=g_{am}\,R^m{}_{bcd}$ (``a manner of speaking: pull the first index down''). This contraction is chart independent, because $g_{am}$ transforms opposite to the upper index of $R^m{}_{bcd}$, ``I need in principle to check that''. It is not the Riemann curvature, which has $a$ up and $b,c,d$ down.
  \item The \textbf{Ricci curvature} is the $(0,2)$ tensor field $R_{ab}:=R^m{}_{amb}$, the contraction of the first and the third index. It could be defined without a metric (``if this is the curvature tensor of some connection I can always do this without a metric''), but on a metric manifold the metric comes in through the back door, because the connection used is the Levi-Civita one.
  \item The \textbf{scalar curvature} (Ricci scalar) is $R:=g^{ab}R_{ab}$: ``this can only be defined on a metric manifold'', one takes the Ricci tensor and contracts both indices by the (inverse) metric; it is truly metric dependent.
  \item The \textbf{Einstein curvature} is the $(0,2)$ tensor field
        \[
          G_{ab}:=R_{ab}-\tfrac12\,g_{ab}\,R,
        \]
        which is the left-hand side of Einstein's equations. All these come from the metric.
\end{enumerate}
\end{definition}
\ts{10}{1}{1:12:12}\ts{10}{1}{1:12:32}\ts{10}{1}{1:12:52}\ts{10}{1}{1:13:13}\ts{10}{1}{1:14:37}\ts{10}{1}{1:14:58}\ts{10}{1}{1:15:18}\ts{10}{1}{1:15:39}\ts{10}{1}{1:16:00}\ts{10}{1}{1:16:20}\ts{10}{1}{1:16:40}\ts{10}{1}{1:17:01}.
There are important properties (the Riemann--Christoffel curvature is antisymmetric in the first two and in the last two indices, and one can exchange the two pairs with nothing happening): ``you will check all of that in the tutorials, very important to practise and very important for the future'' (\cref{sec:t10-1}, exercise 6) \ts{10}{1}{1:15:39}.
\begin{aside}[Einstein's missing term]
The combination $G_{ab}$ plays a role in the Einstein equations ``as the combination that Einstein only realised rather late'' (the name capital $G_{ab}$ ``for whatever reasons''): in the first lecture he said that Einstein wrote $R_{ab}=T_{ab}$ (up to constants) and ``forgot this term'' of $-\tfrac12g_{ab}R$ in 1912. This is the term of which \cref{sec:c09-chart} said that it is missing in the old form of the equations; if $g$ is Lorentzian all this goes through for Lorentzian metrics. ``This combination is nowadays very famous, has very nice properties, we will come to speak of that''
\ts{10}{1}{1:16:00}\ts{10}{1}{1:16:20}\ts{10}{1}{1:16:40}\ts{10}{1}{1:17:01}.
\end{aside}
\begin{coursenote}[Student questions at the end]
\begin{itemize}
  \item \emph{The indices on $g$ in $G$.} A student points at the $g$ with upper indices in the contraction $R=g^{ab}R_{ab}$: yes, these are the components of the so-called inverse metric, a convention; written like that ``it was a mistake according to the notation I chose, but it is a convention'' (see above) \ts{10}{1}{1:17:01}\ts{10}{1}{1:17:24}\ts{10}{1}{1:17:44}\ts{10}{1}{1:18:05}.
  \item \emph{What about the trace of $G$?} One could contract $G_{ab}$ with $g^{ab}$: $g^{ab}R_{ab}-\tfrac12g^{ab}g_{ab}R=R-\tfrac12\,\delta^a{}_aR=\bigl(1-\tfrac d2\bigr)R$ with $d=\dim M$ (4 in four dimensions): in two dimensions the trace vanishes (``that is right, but only in two dimensions''); nothing else happens when one plugs this in \ts{10}{1}{1:18:26}\ts{10}{1}{1:18:46}\ts{10}{1}{1:19:06}\ts{10}{1}{1:19:27}.
\end{itemize}
\end{coursenote}
Schuller invites the audience to the tutorials at 2:30 (``very important exercises that are important for general relativity later on''); one of the three tutorials moved to another room because of the projector \ts{10}{1}{1:19:51}\ts{10}{1}{1:20:11}.

% ==================================================================
\section{Tutorial 9 \& 10}\label{sec:t10-1}

Tutorial 9 \& 10, 1:17:31, \seg{10}{2}{video}, by a tutor, not Schuller. The sheets (Sheet 9 and Sheet 10) are paraphrased, not copied; see \url{https://tales.mbivert.com/on-heraeus-winter-school-tutorials/} (and \url{https://tales.mbivert.com/heraeus/Sheet_09.pdf}, \url{https://tales.mbivert.com/heraeus/Sheet_10.pdf}). (Mis-hearing in the captions: the metric $g$ is often transcribed as ``T'' or ``G''; read $g$.) The sheet starts with true/false questions \ts{10}{2}{0:00}\ts{10}{2}{0:20}.

\subsection*{Exercise 1: true or false}

\paragraph{(A) The Levi-Civita connection.}
\begin{itemize}
  \item It is torsion free: \textbf{true}, only torsion-free connections are considered \ts{10}{2}{0:20}\ts{10}{2}{0:41}.
  \item It is called metric compatible because it satisfies a given equation: \textbf{false} as stated, since the compatibility condition is $\nabla_Xg=0$, the zero tensor field, and the equation on the sheet has a different right-hand side [the exact wording could not be recovered from the captions, ?] \ts{10}{2}{0:41}\ts{10}{2}{1:03}.
  \item It is the unique affine connection that can be constructed from the data of a metric manifold: \textbf{true} (exercise 3) \ts{10}{2}{1:03}.
  \item It gives connection coefficient functions with vanishing symmetric part $\Gamma^a{}_{(bc)}=0$ in any chart: \textbf{false}; this part of the connection coefficients is not tensorial, it may vanish in one chart but surely not in any \ts{10}{2}{1:25}.
  \item It arises by requiring that the auto-parallels be the curves of stationary length: \textbf{true}, seen in the lecture \ts{10}{2}{1:46}.
\end{itemize}
So only statements 1, 3 and 5 are correct \ts{10}{2}{1:46}\ts{10}{2}{2:07}.

\paragraph{(B) Metrics on a $d$-dimensional smooth manifold.}
\begin{itemize}
  \item A metric with signature $(d,d)$ is Riemannian: \textbf{false}; a Riemannian metric has signature $(d,0)$ or $(0,d)$ \ts{10}{2}{2:07}.
  \item A Lorentzian metric has signature $(1,d-1)$: \textbf{true} (or $(d-1,1)$) \ts{10}{2}{2:29}.
  \item Vectors $X\in T_pM$ with $g(X,X)=0$ on a Lorentzian manifold are called null vectors: \textbf{true}, the definition \ts{10}{2}{2:49}.
  \item A metric provides an inner product on each $T_pM$: \textbf{false}; an inner product needs positive definiteness, which holds for Riemannian metrics but for no other signature \ts{10}{2}{2:49}\ts{10}{2}{3:11}.
  \item An inverse metric fed with a covector yields a real number: \textbf{false}; the inverse metric eats two covectors, with one fed in one slot is left, a $(1,0)$ tensor field \ts{10}{2}{3:11}\ts{10}{2}{3:32}.
\end{itemize}

\paragraph{(C) Geodesics and the length of a curve.}
\begin{itemize}
  \item A Riemannian metric gives a well-defined length of a curve: \textbf{true} \ts{10}{2}{3:32}.
  \item The length of a curve is not invariant under reparametrisation: \textbf{false}, it is invariant under increasing bijective reparametrisations (shown below) \ts{10}{2}{3:54}\ts{10}{2}{4:15}.
  \item A geodesic is a curve of minimal length: \textbf{false}; of \emph{stationary} length \ts{10}{2}{4:15}.
  \item A curve is of stationary length iff it solves the Euler--Lagrange equations for $L(x,v)=\sqrt{g(v,v)}$: \textbf{true}, seen in the lecture \ts{10}{2}{4:35}.
  \item A geodesic on a metric manifold satisfies the auto-parallel equation: \textbf{true} for the Levi-Civita connection. Some words must be said: geodesics describe the shortest curves between two points and auto-parallels the straightest curves, two notions that are in general disjoint, but if one chooses the Levi-Civita connection they coincide \ts{10}{2}{4:35}\ts{10}{2}{4:57}\ts{10}{2}{5:18}.
\end{itemize}

\subsection*{Exercise 2: signatures and null vectors}

A symmetric $(0,2)$ tensor $g$ on a three-dimensional real vector space has no eigenvalues but a signature \ts{10}{2}{5:18}\ts{10}{2}{5:38}. Given four sets of null vectors of $g$ (a double cone around a line through the origin, a single point, a straight line through the origin, a plane through the origin), find the possible signatures of $g$ and indicate the vectors of positive length \ts{10}{2}{5:38}\ts{10}{2}{6:00}.
Choose coordinates $x^1,x^2,x^3$ on the vector space.
\begin{enumerate}
  \item \emph{Double cone.} The null vectors satisfy $(x^1)^2+(x^2)^2-(x^3)^2=0$, which is $g(x,x)=0$ for the signature $(+,+,-)$; also $(-,-,+)$ gives the same cone (multiply the equation by $-1$). Positive length (computed directly, ``outside'' for $(+,+,-)$, ``inside'', i.e.\ $|x^3|>\sqrt{(x^1)^2+(x^2)^2}$, for $(-,-,+)$); this is the picture of relevance for general relativity \ts{10}{2}{6:20}\ts{10}{2}{6:42}\ts{10}{2}{7:03}\ts{10}{2}{7:24}\ts{10}{2}{9:54}\ts{10}{2}{10:15}\ts{10}{2}{10:36}.
  \item \emph{A single point (the origin).} $g(x,x)=0$ only for $x=0$: no sum of squares of components with all the same sign can vanish except for the null vector, so $(+,+,+)$ or $(-,-,-)$: positive length for all non-zero vectors in the first case, none in the second \ts{10}{2}{7:24}\ts{10}{2}{7:44}\ts{10}{2}{8:05}\ts{10}{2}{10:36}\ts{10}{2}{10:56}.
  \item \emph{A straight line through the origin.} The null vectors form a line, the third component being free: $(+,+,0)$ or $(-,-,0)$; the first two entries cannot add up to zero otherwise. Positive length: everything outside the line for $(+,+,0)$, none for $(-,-,0)$ \ts{10}{2}{8:27}\ts{10}{2}{8:49}\ts{10}{2}{10:56}\ts{10}{2}{11:16}.
  \item \emph{A plane through the origin.} $(+,0,0)$ or $(-,0,0)$ (the tutor writes the zeros first, ``$0,0,+$''): $x^1,x^2$ may be anything, the single remaining component never adds up to zero; positive length everywhere outside the plane in the first case, none in the second \ts{10}{2}{9:10}\ts{10}{2}{9:31}\ts{10}{2}{11:16}\ts{10}{2}{11:37}.
\end{enumerate}
In general relativity the Lorentzian case, positive vectors inside or outside the cone, is the important one \ts{10}{2}{10:36}.

\subsection*{Exercise 3: the Levi-Civita connection is unique}

For vanishing torsion tensor the Levi-Civita connection is already uniquely determined by the requirement $\nabla g=0$: suppose $\nabla$ is torsion free and metric compatible; $\nabla g$ is a $(0,3)$ tensor field with the derivative slot left open \ts{10}{2}{11:37}\ts{10}{2}{11:58}\ts{10}{2}{12:19}
\begin{enumerate}
  \item \emph{What does $T=0$ imply for the coefficients?} $T^a{}_{bc}=\Gamma^a{}_{bc}-\Gamma^a{}_{cb}=0$, so all antisymmetric components vanish: $\Gamma^a{}_{bc}=\Gamma^a{}_{cb}$ \ts{10}{2}{12:19}\ts{10}{2}{12:39}\ts{10}{2}{13:02}\ts{10}{2}{13:25}.
  \item \emph{Write $\nabla g=0$ in terms of the $\Gamma$'s.} With the abbreviation $g_{bc,a}:=\partial g_{bc}/\partial x^a$:
        \[
          (\nabla_ag)_{bc}=g_{bc,a}-\Gamma^n{}_{ba}\,g_{nc}-\Gamma^n{}_{ca}\,g_{bn}=0,
        \]
        and the two cyclic permutations of $(a,b,c)$, giving three equations \ts{10}{2}{13:25}\ts{10}{2}{13:46}\ts{10}{2}{14:07}\ts{10}{2}{14:28}\ts{10}{2}{14:49}\ts{10}{2}{15:11}\ts{10}{2}{15:31}\ts{10}{2}{15:53}.
  \item \emph{Solve for $\Gamma$.} Add the first two and subtract the third:
        \[
          g_{bc,a}+g_{ca,b}-g_{ab,c}-2\,\Gamma^n{}_{ab}\,g_{cn}=0,
        \]
        because, with the three equations $g_{bc,a}=\Gamma^n{}_{ba}g_{nc}+\Gamma^n{}_{ca}g_{bn}$, $g_{ca,b}=\Gamma^n{}_{cb}g_{na}+\Gamma^n{}_{ab}g_{cn}$, $g_{ab,c}=\Gamma^n{}_{ac}g_{nb}+\Gamma^n{}_{bc}g_{an}$, the symmetry $\Gamma^n{}_{bc}=\Gamma^n{}_{cb}$ and that of the metric make the terms $\Gamma^n{}_{ca}g_{bn}$ and $\Gamma^n{}_{ac}g_{nb}$ cancel, likewise $\Gamma^n{}_{cb}g_{na}$ and $\Gamma^n{}_{bc}g_{an}$, and the two terms $\Gamma^n{}_{ba}g_{nc}=\Gamma^n{}_{ab}g_{cn}$ add up. Multiplying by $g^{cm}$ ($g_{cn}g^{cm}=\delta^m{}_n$):
        \[
          \Gamma^m{}_{ab}=\tfrac12\,g^{cm}\bigl(g_{bc,a}+g_{ca,b}-g_{ab,c}\bigr),
        \]
        exactly the expression of \cref{def:c10-lc}. Hence torsion freeness and $\nabla g=0$ determine the connection coefficients uniquely in terms of the metric \ts{10}{2}{15:53}\ts{10}{2}{16:14}\ts{10}{2}{16:34}\ts{10}{2}{16:55}\ts{10}{2}{17:18}\ts{10}{2}{17:39}\ts{10}{2}{17:59}\ts{10}{2}{18:20}\ts{10}{2}{18:42}\ts{10}{2}{19:04}\ts{10}{2}{19:25}\ts{10}{2}{19:46}\ts{10}{2}{20:08}\ts{10}{2}{20:30}\ts{10}{2}{20:51}\ts{10}{2}{21:12}\ts{10}{2}{21:33}\ts{10}{2}{21:53}\ts{10}{2}{22:14}\ts{10}{2}{22:34}\ts{10}{2}{22:55}\ts{10}{2}{23:15}\ts{10}{2}{23:37}\ts{10}{2}{23:57}.
\end{enumerate}

\subsection*{Exercise 4: massaging the length functional}

\begin{enumerate}
  \item \emph{Reparametrisation invariance} (\cref{thm:c10-reparam}). Let $\tilde\gamma=\gamma\circ\sigma$ with $\sigma$ increasing, smooth and bijective. By definition $L[\tilde\gamma]=\int_0^1\mathrm d\lambda\sqrt{g_{ab}(\tilde\gamma(\lambda))\,\dot{\tilde\gamma}^a(\lambda)\,\dot{\tilde\gamma}^b(\lambda)}$ with $\dot{\tilde\gamma}^a=(x^a\circ\tilde\gamma)'$. By the chain rule $(x^a\circ\gamma\circ\sigma)'(\lambda)=(x^a\circ\gamma)'\bigl(\sigma(\lambda)\bigr)\,\sigma'(\lambda)$, so the factor $\sigma'(\lambda)^2$ comes out of the square root as $\sigma'(\lambda)$, which is positive because $\sigma$ is increasing, and
        \[
          L[\tilde\gamma]=\int_0^1\mathrm d\lambda\,\sigma'(\lambda)\sqrt{g_{ab}\bigl(\gamma(\sigma(\lambda))\bigr)\,\dot\gamma^a(\sigma(\lambda))\,\dot\gamma^b(\sigma(\lambda))}.
        \]
        Substitute $\hat\lambda:=\sigma(\lambda)$, $\mathrm d\hat\lambda=\sigma'(\lambda)\,\mathrm d\lambda$; since $\sigma$ is increasing and bijective the limits stay $0$ and $1$, and one gets exactly the integral for $L[\gamma]$ (with $\hat\lambda$) \ts{10}{2}{24:18}\ts{10}{2}{24:38}\ts{10}{2}{25:00}\ts{10}{2}{25:23}\ts{10}{2}{25:44}\ts{10}{2}{26:06}\ts{10}{2}{26:27}\ts{10}{2}{26:48}\ts{10}{2}{27:08}\ts{10}{2}{27:31}\ts{10}{2}{27:51}\ts{10}{2}{28:12}\ts{10}{2}{28:32}\ts{10}{2}{28:53}\ts{10}{2}{29:14}\ts{10}{2}{29:35}\ts{10}{2}{29:55}\ts{10}{2}{30:18}\ts{10}{2}{30:38}\ts{10}{2}{31:01}\ts{10}{2}{31:24}.
  \item \emph{The square root can be dropped.} Let $T:=g_{ab}\dot q^a\dot q^b$ and $L:=\sqrt T$. The Euler--Lagrange equations for $\sqrt T$ have precisely the same solutions as those for $T$, if for the former one only selects the solutions that satisfy $T=1$ on their parametrisation. Indeed
        \[
          \frac{\mathrm d}{\mathrm d\lambda}\frac{\partial\sqrt T}{\partial\dot q^a}-\frac{\partial\sqrt T}{\partial q^a}=\frac{\mathrm d}{\mathrm d\lambda}\Bigl(\frac1{2\sqrt T}\frac{\partial T}{\partial\dot q^a}\Bigr)-\frac1{2\sqrt T}\frac{\partial T}{\partial q^a}=\frac1{2\sqrt T}\Bigl[\frac{\mathrm d}{\mathrm d\lambda}\frac{\partial T}{\partial\dot q^a}-\frac{\partial T}{\partial q^a}\Bigr]
        \]
        because $T=1$ along the parametrisation, so $\sqrt T$ is a constant that can be pulled out of the derivative with respect to the parameter. The bracket is the Euler--Lagrange expression for $T$. Why bother? Because the square root is awkward and this shows that it does not change the equations, simplifying the computations \ts{10}{2}{31:44}\ts{10}{2}{32:08}\ts{10}{2}{32:32}\ts{10}{2}{32:53}\ts{10}{2}{33:16}\ts{10}{2}{33:37}\ts{10}{2}{33:59}\ts{10}{2}{34:20}\ts{10}{2}{34:40}\ts{10}{2}{35:01}.
\end{enumerate}

\subsection*{Exercise 5: the geodesic equation of the round sphere}

A practical way to quickly determine Christoffel symbols in a concrete case: re-deriving the Euler--Lagrange equations is quicker than using the general formula of exercise 3 \ts{10}{2}{35:01}\ts{10}{2}{35:22}. Derive the geodesic equation for the round sphere of radius $R$ with the metric of \cref{ex:c10-sphere} in the chart $(\theta,\varphi)$. Since $\ddot\gamma^a+\Gamma^a{}_{bc}\dot\gamma^b\dot\gamma^c=0$ with the Levi-Civita $\Gamma$'s is both the geodesic and the auto-parallel equation, one reads the $\Gamma$'s off. Use $L=g_{ab}\dot\gamma^a\dot\gamma^b=R^2\dot\theta^2+R^2\sin^2\theta\,\dot\varphi^2$ (without the root, by exercise 4) \ts{10}{2}{35:22}\ts{10}{2}{35:43}\ts{10}{2}{36:04}\ts{10}{2}{36:25}\ts{10}{2}{36:45}\ts{10}{2}{37:06}\ts{10}{2}{37:26}\ts{10}{2}{37:48}\ts{10}{2}{38:09}\ts{10}{2}{38:30}.
\begin{itemize}
  \item $\theta$: $\dfrac{\mathrm d}{\mathrm d\lambda}\dfrac{\partial L}{\partial\dot\theta}-\dfrac{\partial L}{\partial\theta}=2R^2\ddot\theta-2R^2\sin\theta\cos\theta\,\dot\varphi^2=0$, i.e.
        \[\ddot\theta-\sin\theta\cos\theta\,\dot\varphi^2=0\]
        (dividing by $2R^2$; the tutor writes a stray ``$2R^2$'' in the second term, an error corrected on division) \ts{10}{2}{38:50}\ts{10}{2}{39:13}\ts{10}{2}{39:34}\ts{10}{2}{39:54}\ts{10}{2}{40:15}.
  \item $\varphi$: $\dfrac{\mathrm d}{\mathrm d\lambda}\bigl(2R^2\sin^2\theta\,\dot\varphi\bigr)=2R^2\bigl(2\sin\theta\cos\theta\,\dot\theta\dot\varphi+\sin^2\theta\,\ddot\varphi\bigr)=0$, i.e.
        \[\ddot\varphi+2\cot\theta\,\dot\theta\,\dot\varphi=0\]
        \ts{10}{2}{40:39}\ts{10}{2}{41:01}\ts{10}{2}{41:22}\ts{10}{2}{41:43}\ts{10}{2}{42:04}\ts{10}{2}{42:27}.
\end{itemize}
Comparing with the auto-parallel equation, $\Gamma^1{}_{22}=-\sin\theta\cos\theta$ and, since $\Gamma^2{}_{12}=\Gamma^2{}_{21}$ (torsion free) and the term is $2\Gamma^2{}_{12}\dot\theta\dot\varphi$, $\Gamma^2{}_{12}=\Gamma^2{}_{21}=\cot\theta$ (the tutor first reads off $2\cot\theta$ and corrects by distributing the $2$ over the two symbols) \ts{10}{2}{42:49}\ts{10}{2}{43:10}\ts{10}{2}{43:31}\ts{10}{2}{43:53}\ts{10}{2}{44:13}\ts{10}{2}{44:34}\ts{10}{2}{44:55}.

\subsection*{Exercise 6: algebraic symmetries of the Riemann--Christoffel tensor}

Properties that a plain Riemann curvature does not feature (\cref{def:c10-curv}) \ts{10}{2}{44:55}\ts{10}{2}{45:16}.
\begin{enumerate}
  \item \emph{Derivative of the inverse metric.} Differentiating $g^{ab}g_{bc}=\delta^a{}_c$ (a constant) with respect to $x^d$ and using the product rule gives $\partial_d g^{ab}\,g_{bc}+g^{ab}\,\partial_dg_{bc}=0$, hence $\partial_dg^{ae}=-g^{ab}\,(\partial_dg_{bc})\,g^{ce}$ \ts{10}{2}{45:16}\ts{10}{2}{45:37}\ts{10}{2}{45:58}\ts{10}{2}{46:20}\ts{10}{2}{46:43}\ts{10}{2}{47:03}\ts{10}{2}{47:24}\ts{10}{2}{47:47}\ts{10}{2}{48:09}\ts{10}{2}{48:30}.
  \item \emph{$R_{abcd}$ in normal coordinates.} In normal coordinates at a point $p$ (all $\Gamma$'s vanish there: chapter~\ref{ch:c07}, \cref{thm:c07-normal}; also the first derivatives of the metric vanish at $p$ by $\nabla g=0$), with $g_{ab,cd}:=\partial^2g_{ab}/\partial x^c\partial x^d$,
        \[
          R_{abcd}\big|_p=\tfrac12\bigl(g_{ad,bc}-g_{bd,ac}-g_{ac,bd}+g_{bc,ad}\bigr).
        \]
        One starts from $R_{abcd}=g_{ak}R^k{}_{bcd}$ with $R^k{}_{bcd}=\partial_c\Gamma^k{}_{db}-\partial_d\Gamma^k{}_{cb}+(\Gamma\Gamma)$ (the products of $\Gamma$'s vanish at $p$), inserts the Christoffel symbols, applies the product rule (the terms with a derivative of the inverse metric are proportional to the first derivatives of $g$ and vanish by exercise 6.1) and uses $g_{ak}g^{kl}=\delta_a{}^l$. The tutor makes sign slips on the board (``it is just a calculation with indices, you need to get used to it, it is a matter of concentration''), corrects them and arrives at the formula \ts{10}{2}{48:30}\ts{10}{2}{48:52}\ts{10}{2}{49:12}\ts{10}{2}{49:34}\ts{10}{2}{49:55}\ts{10}{2}{50:16}\ts{10}{2}{50:38}\ts{10}{2}{50:58}\ts{10}{2}{51:20}\ts{10}{2}{51:41}\ts{10}{2}{52:02}\ts{10}{2}{52:25}\ts{10}{2}{52:48}\ts{10}{2}{53:10}\ts{10}{2}{53:31}\ts{10}{2}{53:52}\ts{10}{2}{54:16}\ts{10}{2}{54:38}\ts{10}{2}{54:59}\ts{10}{2}{55:21}\ts{10}{2}{55:43}\ts{10}{2}{56:04}\ts{10}{2}{56:24}\ts{10}{2}{56:44}\ts{10}{2}{57:06}\ts{10}{2}{57:28}\ts{10}{2}{57:50}\ts{10}{2}{58:12}\ts{10}{2}{58:32}\ts{10}{2}{58:53}\ts{10}{2}{59:13}\ts{10}{2}{59:34}\ts{10}{2}{59:55}\ts{10}{2}{1:00:16}\ts{10}{2}{1:00:43}\ts{10}{2}{1:01:04}\ts{10}{2}{1:01:25}\ts{10}{2}{1:01:45}.
  \item \emph{The symmetries} (the formula above, now used in normal coordinates, holds for the tensor components in general since they are tensor equations): $R_{abcd}=-R_{bacd}$ (exchange of the first two indices), $R_{abcd}=-R_{abdc}$ (the last two), $R_{abcd}=R_{cdab}$ (exchange of the two pairs), and the cyclic identity $R_{a[bcd]}=0$. Each is checked by inserting the formula and comparing the four terms: they cancel pairwise or map into each other, using the symmetry of the metric and Schwarz's rule (``since we derive two times in the same coordinate system''). The full verification is in the appendix (\cref{app:c10-riemann}) \ts{10}{2}{1:02:06}\ts{10}{2}{1:02:31}\ts{10}{2}{1:03:55}\ts{10}{2}{1:07:07}\ts{10}{2}{1:17:13}.
\end{enumerate}

% ==================================================================
\section*{Key concepts}
\begin{keyconcepts}
\textbf{Retrieval questions.}
\begin{enumerate}
  \item Define a metric on a smooth manifold: which two conditions? What is the musical map, and why must it be invertible?
  \item What is the so-called inverse metric and why ``so-called''? Why does the lecture never suppress the metric when raising or lowering?
  \item What is the signature of a symmetric $(0,2)$ tensor, and why is it not a list of eigenvalues? Define Riemannian, Lorentzian and pseudo-Riemannian.
  \item Define speed and length of a curve from a metric. Compute the length of the equator of the round sphere with a non-uniform parametrisation. Why is the length invariant under increasing reparametrisation?
  \item What is a geodesic (stationary, not minimal)? Derive the geodesic equation from the Euler--Lagrange equations of $\sqrt{g(\dot x,\dot x)}$, choosing unit speed.
  \item What are the Christoffel symbols and the Levi-Civita connection? Which two conditions determine it uniquely?
  \item Define the Riemann--Christoffel, Ricci, scalar and Einstein curvatures. What is the trace of $G_{ab}$ in $d$ dimensions?
\end{enumerate}
\textbf{Concepts.}
\begin{itemize}
  \item Metric $g$ symmetric and non-degenerate, musical map $\flat$, $g^{-1}$ (\cref{def:c10-metric}, \cref{def:c10-inverse}); round sphere (\cref{ex:c10-sphere}).
  \item Signature $(p,q)$; Riemannian, Lorentzian, pseudo-Riemannian (\cref{def:c10-signature}).
  \item Speed, length, reparametrisation invariance (\cref{def:c10-speed}, \cref{thm:c10-reparam}); geodesics as stationary curves, Euler--Lagrange (\cref{def:c10-geodesic}, \cref{thm:c10-el}).
  \item Christoffel symbols $\Gamma^q{}_{ij}=\tfrac12g^{qm}(\partial_ig_{mj}+\partial_jg_{mi}-\partial_mg_{ij})$ and the Levi-Civita connection (\cref{def:c10-lc}); metric compatibility $\nabla g=0$ plus zero torsion.
  \item $R_{abcd}$, $R_{ab}$, $R$ and $G_{ab}=R_{ab}-\tfrac12g_{ab}R$ (\cref{def:c10-curv}).
\end{itemize}
\end{keyconcepts}

\flushside
\end{paracol}
